\subsection{结构水热共同参数设计}\label{app:parameter-table}
下表给出空间填充设计的19项参数上下限及取样尺度。四机制共享256个参数点，两强度政策形成2048个设计情景；参数设计与概率先验分别登记。实际各参数点见相应参数表。
\begingroup\small\begin{longtable}{p{3.8cm}p{2.2cm}rrp{1.6cm}}
\caption{共同参数设计范围。}\label{tab:app-parameters}\\
\toprule 参数 & 单位 & 下限 & 上限 & 取样尺度\\\midrule\endfirsthead
\toprule 参数 & 单位 & 下限 & 上限 & 取样尺度\\\midrule\endhead
结构面倾角 & $^\circ$ & 30 & 60 & 线性\\
法向厚度 & m & 10 & 60 & 线性\\
岩体体积分数 & 1 & 0.2 & 0.8 & 线性\\
岩桥承载面积比例 & 1 & 0 & 0.25 & 线性\\
岩桥有效黏聚力 & kPa & 300 & 3000 & 对数\\
岩面摩擦系数 & 1 & 0.45 & 0.85 & 线性\\
初始裂隙温度 & $^\circ$C & -6 & -0.8 & 线性\\
热扩散率 & m$^2$/s & 9e-07 & 2.1e-06 & 线性\\
热路径深度 & m & 0.5 & 20 & 对数\\
站点至源区温度偏移 & $^\circ$C & -18 & -6 & 线性\\
储水量/水头比 & 1 & 0.002 & 0.08 & 对数\\
排水时间常数 & h & 3 & 240 & 对数\\
最大连通水头 & m & 2 & 40 & 对数\\
有效入渗倍率 & 1 & 0.25 & 3 & 对数\\
度日融化系数 & mm/($^\circ$C$\cdot$d) & 0.5 & 6 & 线性\\
损伤速率系数 & h$^{-1}$ & 1e-07 & 0.0002 & 对数\\
应力比幂次 & 1 & 2 & 6 & 线性\\
初始损伤 & 1 & 0.05 & 0.7 & 线性\\
损伤排水耦合系数 & 1 & 0 & 3 & 线性\\
\bottomrule\end{longtable}\endgroup
\subsection{解析热机制分类与裕度补充}
分类按相应初态与首次过阈定义，在同一设计表内汇总；初始越阈与未来越阈分列。
\begin{table}[htbp]\centering\small\caption{最终解析热版本的机制分类与最小稳定裕度中位数。}\label{tab:app-mechanisms}
\begin{tabular}{llrrrr}\toprule 机制 & 强度政策 & 未来越阈 & 初始越阈 & 时域内保持 & 裕度中位数\\\midrule
S & 线性延伸 & 4 & 31 & 221 & 0.6901\\
S+T & 线性延伸 & 28 & 31 & 197 & 0.4853\\
S+H & 线性延伸 & 6 & 31 & 219 & 0.6623\\
S+T+H & 线性延伸 & 31 & 31 & 194 & 0.4718\\
S & 暖端岩接触 & 4 & 31 & 221 & 0.6901\\
S+T & 暖端岩接触 & 27 & 31 & 198 & 0.4853\\
S+H & 暖端岩接触 & 6 & 31 & 219 & 0.6623\\
S+T+H & 暖端岩接触 & 31 & 31 & 194 & 0.4718\\
\bottomrule\end{tabular}\end{table}
